\documentclass[11pt,a4,oneside]{article}
\usepackage{graphicx} % Required for inserting images
\usepackage{geometry}
\usepackage{longtable}
\usepackage{multirow}
\usepackage{caption}
\usepackage{subcaption}
\usepackage[table,xcdraw]{xcolor}
\usepackage{bibentry}
\usepackage[skip=10pt plus1pt, indent=40pt]{parskip}
\usepackage[portuguese]{babel}

\geometry{
 a4paper,
left=3cm,
top=3cm,
right=2cm,
bottom=2cm
 }
\usepackage{fancyhdr}
\usepackage{biblatex} %Imports biblatex package
\addbibresource{references.bib} %Import the bibliography file
\linespread{1.2}

\title{\textit{Checklist} para verificação de aderência do tema de Trabalho de Conclusão de Curso I ao curso de Ciência da Computação.}
\date{Fevereiro 2026}

\begin{document}
\pagestyle{fancy}
%\fancyhead[C]{In the centre of the header on all pages}
\pagestyle{fancy}
\fancyhead{}
\fancyhead[RO,LE]{\textbf{Checklist  - TCC I}}

\maketitle
\begin{itemize}
\item \textbf{Modalidade:} Acadêmico - Projeto - Grupo
\item \textbf{Tema:} Identificação de agentes estressores no Metrô de São Paulo utilizando algoritmos de aprendizado não-supervisionado
\item \textbf{Alunos:}
\begin{itemize}
    \item Lucas Kerr do Amaral (RA: 22.123.032-9)
    \item Marcela Nalesso (RA: 22.222.011-3)
\end{itemize}
\item \textbf{Orientador:} Prof. Dr. Leonardo Anjoletto Ferreira
\item \textbf{Resumo:} Os funcionários do Metrô reportaram que gostariam de compreender melhor como é possível melhorar a qualidade da saúde dos funcionários do Metrô e, para isso uma possibilidade é aplicar questionários ou desenvolver aplicativos que monitorem a saúde mental destes funcionários e com inteligência artificial descobrir padrões que podem ser úteis, por exemplo: mulheres até 40 anos e homens de mais de 50 que trabalham no setor de atendimento a usuários na catraca, no período matutino, possuem indicativos de estarem sofrendo \textit{burnout}. Estamos pesquisando se existem questionários prontos para isso e quais algoritmos de inteligência artificial podem nos ajudar nessa tarefa.
\end{itemize}
\section*{Lista de disciplinas aderentes}
\begin{itemize}
\item \textbf{CS1711 Comunicação e Expressão:} Utilizaremos conceitos de apresentação e escritas formais para .......
\item \textbf{CS1711 Desenvolvimento WEB:} Se contruirmos um site, acho válido colocar os conceitos de html, js e css (caso usarmos)
\item \textbf{CC1612 Fundamentos de Algoritmos:} Utilizaremos conceitos de lógica como condições e leitura de arquivos vistos na disciplina para  auxiliar na implementação do projeto.
\item \textbf{CC5232 Banco de Dados:} A partir dos questionários realizados talvez seja necessário implementar um Banco de Dados para armazenar esses dados e ter um histórico caso o Metrô realize ações afirmativas entre seus funcionários e queria comparar resultados.
\item \textbf{CA5411 Probabilidade e Estatística:} Pelo que temos lido, a área de IA e estatística são muito relacionadas. Além disso a orientadora comentou que será necessário estudar métodos de normalização dos dados, o que é um conceito vindo dessa área.
\item \textbf{CC7711 Inteligência Artificial e Robótica:} Como vamos utilizar métodos de IA, é bem possível que utilizaremos conceitos dessa disciplina. Questionamos o professor da disciplina sobre o nosso problema e ele indicou que no nosso caso, precisaremos nos aprofundar em conceitos relacionados a algoritmos não supervisionados.
\end{itemize}
\end{document}
